\exercisegroup

\begin{enumerate}
\def\labelenumi{\arabic{enumi})}
\tightlist
\item
  Let \((\mathrm{Q}_{i})_{i\in\mathrm{I}}\) be a family of quivers. We denote by S the sum set of the family \((\mathrm{Som}(\mathrm{Q}_{i}))\), by F the sum set of the family \((\mathrm{Fl}(\mathrm{Q}_{i}))\), and by o and t the maps from F to S induced by the origin and target maps of the quivers \(Q_{i}\).
\end{enumerate}

\begin{enumerate}
\def\labelenumi{\alph{enumi})}
\item
  Prove that \(\mathrm{Q} = (\mathrm{S}, \mathrm{F}, o, t)\) is a quiver. We call it the sum quiver of the family \((\mathrm{Q}_{i})\).
\item
  For \(i \in I\) there exists a unique quiver morphism \(j_{i}\) from \(Q_{i}\) to Q such that the map \(\operatorname{Som}(j_{i})\) is the canonical map from \(\operatorname{Som}(Q_{i})\) into S and the map \(\operatorname{Fl}(j_{i})\) is the canonical map from \(\operatorname{Fl}(Q_{i})\) into S.
\item
  Let K be a quiver and, for every \(i \in I\), let \(\varphi_{i}\) be a quiver morphism from \(Q_{i}\) to K. Show that there exists a unique quiver morphism \(\varphi\) from Q to K such that one has \(\varphi \circ j_{i} = \varphi_{i}\) for all \(i \in I\).
\end{enumerate}

\begin{enumerate}
\def\labelenumi{\arabic{enumi})}
\setcounter{enumi}{1}
\tightlist
\item
  Let \(Q = (S, F, o, t)\) be a quiver. For every \(s \in S\), we call \(d_{+}(s)\) (resp. \(d_{-}(s)\)) the outgoing (resp. incoming) degree of Q at s and define \(d_{+}(s)\) (resp. \(d_{-}(s)\)) as the cardinality of the set \(\{f \in F \mid o(f) = s\}\) (resp. \(\{f \in F \mid t(f) = s\}\)). We say that Q is locally finite if \(d_{+}(s)\) and \(d_{-}(s)\) are finite for every \(s \in S\), and that it is finite if the sets F and S are finite.
\end{enumerate}

Prove the equalities. \(\sum_{s\in S}d_{+}(s) = \sum_{s\in S}d_{-}(s) = \mathrm{Card}(\mathrm{F})\)

\begin{enumerate}
\def\labelenumi{\arabic{enumi})}
\setcounter{enumi}{2}
\tightlist
\item
  Let \(\mathrm{Q} = (\mathrm{S},\mathrm{F},o,t)\) be a locally finite quiver. We define the matrix \(M_{\mathrm{Q}} = (m_{i,j})\) of type (S,S) with elements in \(\mathbf{Z}\), by
\end{enumerate}

\[
 m _ {i, j} = \operatorname{Card} \left\{f \in \mathrm{F} \mid o (f) = i \text {and} t (f) = j \right\}.
\]

The matrix \(M_{\mathrm{Q}}\) is called the adjacency matrix of Q.

\begin{enumerate}
\def\labelenumi{\alph{enumi})}
\item
  Given quivers Q and Q', give a necessary and sufficient condition on their adjacency matrices for Q and Q' to be isomorphic.
\item
  Express in terms of the adjacency matrix \(M_{Q}\) the number of paths of length n in Q whose origin and term are given.
\item
  Prove that the adjacency matrix of the underlying quiver of a graph is symmetric.
\end{enumerate}

\begin{enumerate}
\def\labelenumi{\arabic{enumi})}
\setcounter{enumi}{3}
\tightlist
\item
  Let \(Q = (S, F, o, t)\) be a quiver. We denote by I the sum set of the family \((S, F)\); we identify S and F with subsets of I. Let k be a commutative ring. The k-algebra of the quiver Q is called the quotient \(A_{Q}\) of the free associative k-algebra on the set I by the two-sided ideal generated by the elements \(X_{s}^{2} - X_{s}\) (for \(s \in S\)), \(X_{u}X_{v} - X_{v}X_{u}\) (for \(u, v \in S\)), \(X_{f}X_{t(f)} - X_{f}\) and \(X_{o(f)}X_{f} - X_{f}\) (for \(f \in F\)). For \(i \in I\), denote by \(x_{i}\) the image of \(X_{i}\) in the algebra \(A_{Q}\).
\end{enumerate}

\begin{enumerate}
\def\labelenumi{\alph{enumi})}
\item
  Show that one has \(\sum_{s\in S}x_s = 1\)
\item
  Prove that we have \(x_{f}x_{g}=0\) if f and g are arrows of Q that are not composable.
\item
  For every path \(c = (a_{0}, f_{1}, \ldots, f_{n}, a_{n})\) in Q, we set \(x_{c} = x_{a_{0}} x_{f_{1}} \ldots x_{f_{n}}\). Let c and \(c'\) be paths in Q; show that we have \(x_{c} x_{c'} = x_{cc'}\) if c and \(c'\) are composable, and \(x_{c} = x_{c'}\) otherwise.
\item
  Let C be the set of paths in the quiver Q; prove that the family \((x_{c})_{c\in C}\) is a basis of the k-module \(A_{Q}\).
\item
  Let \(J_{Q}\) be the two-sided ideal of \(A_{Q}\) generated by the elements of the form \(x_{f}\), for \(f \in F\). Prove that the algebra \(A_{Q}/J_{Q}\) is isomorphic to \(k^{(S)}\).
\item
  For the k-algebra \(A_{Q}\) to be a finitely generated k-module, it is necessary and sufficient that the set of vertices of Q be finite and that every loop in Q have length zero. The ideal \(J_{Q}\) is then nilpotent.
\end{enumerate}

We retain these notations for the rest of the exercises of §1. When we speak of an \(A_{Q}\)-module (without further specification), it is a right module. If A is a ring, an A-module M (without further specification) will be a right A-module; if a is an element of A, we denote by \(a_{M}\) or \((a)_{M}\) the homothety \(x \mapsto xa\) of M.

\begin{enumerate}
\def\labelenumi{\arabic{enumi})}
\setcounter{enumi}{4}
\tightlist
\item
  \begin{enumerate}
  \def\labelenumii{\alph{enumii})}
  \tightlist
  \item
    Prove that the quadruple \(Q^{\circ} = (S, F, t, o)\) is a quiver. It is called the opposite quiver of Q.
  \end{enumerate}
\end{enumerate}

\begin{enumerate}
\def\labelenumi{\alph{enumi})}
\setcounter{enumi}{1}
\tightlist
\item
  The algebra \(A_{Q^{\circ}}\) is isomorphic to the opposite algebra of the algebra \(A_{Q}\) .
\end{enumerate}

\begin{enumerate}
\def\labelenumi{\arabic{enumi})}
\setcounter{enumi}{5}
\tightlist
\item
  \begin{enumerate}
  \def\labelenumii{\alph{enumii})}
  \tightlist
  \item
    Let \((M_s)_{s\in\operatorname{Som}(Q)}\) be a family of \(k\)-modules and let \((u_f)_{f\in\operatorname{Fl}(Q)}\) be a family where, for each \(f\in\operatorname{Fl}(Q)\), \(u_f\) is a linear map from \(M_{o(f)}\) to \(M_{t(f)}\). Let M be the direct sum of the family \((M_s)\); for \(s\in S\), denote by \(j_s\) the canonical injection of \(M_s\) into M and \(p_s\) the canonical projection of M onto \(M_s\). Prove that there exists a unique structure of \(A_Q\)-module on M such that one has \((x_s)_M = j_s \circ p_s\) and \((x_f)_M = j_{t(f)} \circ u_f \circ p_{o(f)}\) for all \(s \in \operatorname{Som}(Q)\) and every \(f \in \operatorname{Fl}(Q)\).
  \end{enumerate}
\end{enumerate}

Conversely, every \(A_{Q}\) -module is of this form.

\begin{enumerate}
\def\labelenumi{\alph{enumi})}
\setcounter{enumi}{1}
\tightlist
\item
  If \(\Lambda\) is a \(k\)-module and \(s \in \operatorname{Som}(\mathbf{Q})\), denote by \(\Lambda(s)\) the \(\mathrm{A}_{\mathrm{Q}}\)-module associated with the family \((\mathrm{M}_i)_{i \in \operatorname{Som}(\mathbf{Q})}\) of \(k\)-modules given by \(\mathrm{M}_i = \Lambda\) if \(i = s\) and \(\mathrm{M}_i = 0\) otherwise, and with the family \((u_f)_{f \in \operatorname{Fl}(\mathbf{Q})}\) of linear maps, where \(u_f = 0\) for all \(f \in \operatorname{Fl}(\mathbf{Q})\).
\end{enumerate}

Let \(\Lambda\) and \(\Lambda'\) be nonzero \(k\)-modules. Let \(s\) and \(s'\) be vertices of Q. Prove that \(\Lambda(s)\) is isomorphic to \(\Lambda'(s')\) if and only if \(\Lambda\) and \(\Lambda'\) are isomorphic as \(k\)-modules and \(s = s'\).

\begin{enumerate}
\def\labelenumi{\alph{enumi})}
\setcounter{enumi}{2}
\item
  If \(\Lambda\) is a simple \(k\)-module, then \(\Lambda(s)\) is a simple \(\mathrm{A}_{\mathrm{Q}}\)-module for every vertex \(s\).
\item
  Assume further that every loop in Q has length zero. Prove that every simple \(A_Q\)-module is of the form \(\Lambda(s)\), where \(\Lambda\) is a simple \(k\)-module and s is a vertex of Q.
\end{enumerate}

\begin{enumerate}
\def\labelenumi{\arabic{enumi})}
\setcounter{enumi}{6}
\tightlist
\item
  Assume that \(k\) is an algebraically closed field. For every \(s \in \operatorname{Som}(Q)\), we set \(\mathrm{P}(s) = x_s\mathrm{A}_{\mathrm{Q}}\).
\end{enumerate}

\begin{enumerate}
\def\labelenumi{\alph{enumi})}
\item
  Prove that \(\mathrm{P}(s)\) is a projective and indecomposable \(A_{Q}\)-module, and that its socle is isomorphic to the \(A_{Q}\)-module \(k(s)\). Prove also that \(\mathrm{P}(s)/\mathrm{P}(s)\mathrm{J}_{\mathrm{Q}}\) is isomorphic to \(k(s)\).
\item
  Assume that every loop in Q has length zero. Prove that, for every \(s \in \operatorname{Som}(Q)\), the canonical homomorphism from \(k\) to \(\operatorname{End}(\mathrm{P}(s))\) is an isomorphism. Prove that, for every projective and indecomposable \(A_Q\)-module P, there exists a unique vertex s of Q such that P is isomorphic to \(\mathrm{P}(s)\).
\end{enumerate}

\begin{enumerate}
\def\labelenumi{\arabic{enumi})}
\setcounter{enumi}{7}
\tightlist
\item
  For \(s \in \operatorname{Som}(Q)\) we set \(\mathrm{P}(s) = x_s\mathrm{A}_{\mathrm{Q}}\) . Let M be a \(\mathrm{A}_{\mathrm{Q}}\) -module.
\end{enumerate}

\begin{enumerate}
\def\labelenumi{\alph{enumi})}
\tightlist
\item
  Let \(f \in \mathrm{Fl}(\mathbf{Q})\), let \(u = o(f)\) and \(v = t(f)\). Prove that there exists a unique pair \((\varphi_f', \varphi_f'')\) of homomorphisms of \(\mathrm{A}_{\mathrm{Q}}\)-modules
\end{enumerate}

\[
\varphi_ {f} ^ {\prime} \colon \mathrm{M} x _ {u} \otimes_ {k} \mathrm{P} (v) \to \mathrm{M} x _ {u} \otimes_ {k} \mathrm{P} (u), \quad \varphi_ {f} ^ {\prime \prime} \colon \mathrm{M} x _ {u} \otimes_ {k} \mathrm{P} (v) \to \mathrm{M} x _ {v} \otimes_ {k} \mathrm{P} (v),
\]

 such that \(\varphi_f'(m \otimes p) = m \otimes x_f p\) and \(\varphi_f''(m \otimes p) = mx_f \otimes p\) for all \(m \in \mathrm{M}x_u\) and every \(p \in \mathrm{P}(v)\).

\begin{enumerate}
\def\labelenumi{\alph{enumi})}
\setcounter{enumi}{1}
\item
  Let \(s \in \text{Som}(Q)\). Prove that there exists a unique homomorphism of \(A_{Q}\)-modules \(\psi_{s}: Mx_{s} \otimes_{k} P(s) \to M\) such that \(\psi_{s}(m \otimes p) = mp\) for \(p \in P(s)\) and \(m \in Mx_{s}\).
\item
  Set \(\varphi = \bigoplus \varphi_f' - \bigoplus \varphi_f''\) and \(\psi = \sum \psi_s\). Show that one has an exact sequence of \(\mathrm{A}_{\mathrm{Q}}\)-modules
\end{enumerate}

\[
0 \to \bigoplus_ {f \in \operatorname{Fl} (\mathrm{Q})} \operatorname{M} x _ {o (f)} \otimes_ {k} \operatorname{P} (t (f)) \xrightarrow {\varphi} \bigoplus_ {s \in \operatorname{Som} (\mathrm{Q})} \operatorname{M} x _ {s} \otimes_ {k} \operatorname{P} (s) \xrightarrow {\psi} \operatorname{M} \to 0.
\]

\(d\)) Show that every right ideal of the algebra \(\mathrm{A_Q}\) is projective.

\begin{enumerate}
\def\labelenumi{\arabic{enumi})}
\setcounter{enumi}{8}
\tightlist
\item
  \begin{enumerate}
  \def\labelenumii{\alph{enumii})}
  \tightlist
  \item
    Let M and N be \(A_Q\)-modules; for \(s \in \operatorname{Som}(Q)\), set \(M_s=Mx_s\) and \(N_s=Nx_s\). For \(\varphi = (\varphi_s)_{s \in \operatorname{Som}(Q)}\) and \(f \in \operatorname{Fl}(Q)\), define a map \(\theta_f(\varphi): Mx_{o(f)} \to Nx_{t(f)}\) by setting
  \end{enumerate}
\end{enumerate}

\[
\theta_ {f} (\varphi) (m) = \varphi_ {o (f)} (m) x _ {f} - \varphi_ {t (f)} (m x _ {f}).
\]

Prove that the map

\[
\theta \colon \prod_ {s \in \operatorname{Som} (\mathrm{Q})} \operatorname{Hom} _ {k} \left(\mathrm{M} _ {s}, \mathrm{N} _ {s}\right)\rightarrow \prod_ {f \in \operatorname{Fl} (\mathrm{Q})} \operatorname{Hom} \left(\mathrm{M} _ {o (f)}, \mathrm{N} _ {t (f)}\right)
\]

 given by \(\varphi\mapsto(\theta_{f}(\varphi))_{f\in\mathrm{Fl}(Q)}\) is linear.

\begin{enumerate}
\def\labelenumi{\alph{enumi})}
\setcounter{enumi}{1}
\item
  Prove that \(\operatorname{Ker}(\theta)\) is isomorphic to \(\operatorname{Hom}_{\mathrm{A}_{\mathrm{Q}}}(\mathrm{M},\mathrm{N})\) and that \(\operatorname{Coker}(\theta)\) is isomorphic to \(\operatorname{Ext}_{\mathrm{A}_{\mathrm{Q}}}^{1}(\mathrm{M},\mathrm{N})\). Prove that \(\operatorname{Ext}_{\mathrm{A}_{\mathrm{Q}}}^{p}(\mathrm{M},\mathrm{N}) = 0\) for every integer \(p \geqslant 2\).
\item
  Suppose that k is a commutative field and that the quiver Q is finite. If M and N are finite-dimensional k-vector spaces, then \(\operatorname{Ext}_{\mathrm{A}_{\mathrm{Q}}}^{p}(\mathrm{M},\mathrm{N})\) is a finite-dimensional k-vector space for every integer \(p \geqslant 0\), is zero for \(p \geqslant 2\), and
\end{enumerate}

\[
\begin{array}{l} \dim_ {k} (\mathrm{Hom} _ {\mathrm{A} _ {\mathrm{Q}}} (\mathrm{M}, \mathrm{N})) - \dim_ {k} (\mathrm{Ext} _ {\mathrm{A} _ {\mathrm{Q}}} ^ {1} (\mathrm{M}, \mathrm{N})) \\ = \sum_ {s \in \mathrm{Som} (\mathrm{Q})} \dim (\mathrm{M} _ {s}) \dim (\mathrm{N} _ {s}) - \sum_ {f \in \mathrm{Fl} (\mathrm{Q})} \dim (\mathrm{M} _ {o (f)}) \dim (\mathrm{N} _ {t (f)}). \end{array}
\]

\(d\)) Assume that \(k\) is a commutative field. Prove that for every pair \((u,v)\) of vertices of Q, the dimension of the vector space \(\mathrm{Ext}_{\mathrm{A_Q}}^1 (k(u),k(v))\) is equal to the cardinality of the set of arrows of Q with origin \(v\) and target \(u\).

\begin{enumerate}
\def\labelenumi{\arabic{enumi})}
\setcounter{enumi}{9}
\item
  We assume that \(k\) is an algebraically closed field. Let Q and Q' be finite quivers in which every loop has length zero. Prove that the algebra A \(_{Q'}\) is Morita equivalent (A, VIII, §6) to the algebra A \(_{Q}\) if and only if the quivers Q and Q' are isomorphic. (Take as modules M, N the \(k\) -simple modules \(k(s)\) , for \(s \in \text{Som}(Q)\) .)
\end{enumerate}

